\chapter{String Matching}

In Chapter~30, we explored the algebraic structure of modular arithmetic, prime numbers, and public-key cryptography. Now we pivot from discrete arithmetic on numbers to the processing of text.

Open any document on your computer and press the keys that mean ``find.'' Type a
word, and every occurrence lights up at once. You have just asked a machine one
of the oldest questions in text processing: \emph{where, inside this long string,
does that short string appear?} It feels instant, but underneath there is an
algorithm, and the obvious way of doing it is not the way your editor actually
does it. The obvious way can be made to crawl. The better ways -- there are
several -- never crawl, and that gap is what this chapter is about.

We begin with the naive sliding-window approach, examine how Rabin--Karp converts substring comparisons into rolling modular hashes, construct deterministic finite automata that track match progress, and culminate in the Knuth--Morris--Pratt algorithm, whose prefix function achieves optimal linear time without ever backtracking in the text.

\chapterline

\section{Setting the stage}
\label{ch31:setup}

\begin{historybox}
In the late 1960s and early 1970s, computing was escaping the rigid world of batch-processed punched cards and entering the era of interactive time-sharing terminals (such as the DEC PDP-10 and SDS 940). Programmers were suddenly writing and editing large source files directly on screen using early text editors like \texttt{QED} and \texttt{TECO}. They needed a way to search for strings in real time. But the physical machinery was unforgiving: text was stored on magnetic core memory, rotating drums, or sequential magnetic tape. Backtracking---rewinding the tape or sliding an internal buffer pointer backward to re-read characters---was computationally painful and mechanically brutal.

The naive search algorithm, which simply resets and shifts by one character whenever a match fails, suffered from a catastrophic blind spot: on repetitive text like searching for $\texttt{a}^{m-1}\texttt{b}$ inside $\texttt{a}^n$, it repeatedly scans $m-1$ matching characters only to fail on the last letter, rewinding the pointer and restarting from scratch. This quadratic crawl of $\Theta(n m)$ operations could stall an interactive time-shared mainframe.

The breakthrough arrived through an unexpected collision of theory and practical hacking. In 1970, James H. Morris was writing a text editor at UC Berkeley and grew exasperated by this rewinding behavior; he devised a clever way to keep the text pointer moving forward by exploiting the pattern's internal structure. Meanwhile at Stanford, Donald Knuth was studying a theoretical result newly proved by Stephen Cook (1970): Cook had demonstrated that any two-way deterministic pushdown automaton could be simulated on a random-access machine in linear time. When Knuth's student Vaughan Pratt observed that string matching was precisely a two-way automaton problem, Knuth set out to turn Cook's abstract theorem into an explicit, highly practical algorithm. Knuth distilled the transition table into an elegant $m$-element integer array---the \emph{prefix function}---and proved that exact matching could be performed in strictly linear $\Theta(n)$ time with \emph{zero} backward movement of the text pointer. Knuth, Morris, and Pratt joined forces to publish their unified algorithm in 1977.

That same year, Robert S. Boyer and J. Strother Moore turned the entire problem upside down. They asked: why match characters from left to right? By aligning the pattern and checking characters \emph{from right to left}, a single mismatch against a character that does not appear anywhere in the pattern allows the entire pattern window to skip forward by its full length $m$ in one bound. This unlocked \emph{sublinear} average-case search times of $O(n/m)$ on natural language texts. A decade later, in 1987, Michael O. Rabin and Richard M. Karp introduced an algebraic alternative: the rolling fingerprint hash, treating substrings as base-$d$ polynomials modulo a prime $q$, computing window updates in $O(1)$ arithmetic steps. While exact single-pattern matching is solved in optimal linear time, a major open frontier is compressed-index pattern matching on massive, highly repetitive genomic collections: querying petabyte-scale pan-genomes directly in compressed space without decompression, using Wheeler graphs and $r$-index data structures.
\end{historybox}

Let us fix notation once and carefully, because string matching is a place where
sloppy indexing causes real pain.

We are given a \emph{text} $T$, an array of $n$ characters $T[1\,..\,n]$, and a
\emph{pattern} $P$, an array of $m$ characters $P[1\,..\,m]$, with $m \le n$.
The characters are drawn from a finite \emph{alphabet} $\Sigma$. For English
text $\Sigma$ might be the $26$ letters plus punctuation and a space; for DNA it
is just $\{\texttt{A},\texttt{C},\texttt{G},\texttt{T}\}$; for raw bytes it is
the $256$ values $0$ through $255$. We write $|\Sigma|$ for the alphabet size.

\begin{definition}[Valid shift]
\label{ch31:def-shift}
We say the pattern $P$ \emph{occurs with shift} $s$ in the text $T$ if
$0 \le s \le n - m$ and
\[
  T[s+1\,..\,s+m] = P[1\,..\,m],
\]
that is, $T[s+j] = P[j]$ for every $j = 1, 2, \ldots, m$. When this holds we
call $s$ a \emph{valid shift}; otherwise $s$ is an \emph{invalid shift}. The
\emph{string-matching problem} is to find all valid shifts.
\end{definition}

The shift $s$ is best read as ``how far the pattern has slid to the right before
it lines up.'' Shift $0$ means the pattern sits at the very front of the text. If
$\texttt{abcab}$ is the text and $\texttt{ab}$ is the pattern, then $s=0$ and
$s=3$ are valid, because the pattern lines up at positions $1$ and $4$.

\begin{intuitionbox}
Picture two strips of paper. The long one is the text, laid flat on the table.
The short one is the pattern, sitting on top of it. You slide the pattern strip
to the right, one notch at a time, and at each resting place you peer down
through the layers and check: do all $m$ characters underneath match the $m$
characters on top? Each resting place is a \emph{shift}. There are exactly
$n-m+1$ of them, from $s=0$ (flush left) to $s=n-m$ (flush right). String
matching is the discipline of checking those resting places \emph{without doing
more peering than you have to}.
\end{intuitionbox}

Figure~\ref{ch31:fig-shifts} draws the strips for a tiny example: the pattern
sliding across the text, one shift at a time.

\begin{figure}[H]
\centering
\begin{tikzpicture}[
  font=\ttfamily\small,
  cell/.style={draw,minimum size=6mm,inner sep=0pt},
  tcell/.style={cell,fill=sourceblue!12},
  pcell/.style={cell,fill=quotientorange!22},
  good/.style={cell,fill=targetgreen!28},
  lbl/.style={font=\sffamily\scriptsize}
]
% text row: abcab
\foreach \c [count=\i from 0] in {a,b,c,a,b} {
  \node[tcell] (t\i) at (\i*0.65,0) {\c};
}
\node[lbl,left=2mm of t0] {$T$};
% shift 0: pattern ab over positions 0,1 -> a,b match -> valid
\node[lbl,left=2mm of t0,yshift=-7mm] {$s=0$};
\node[good] at (0*0.65,-0.7) {a};
\node[good] at (1*0.65,-0.7) {b};
\node[lbl,right=8mm of t4,yshift=-7mm] {valid};
% shift 1: positions 1,2 -> b,c vs a,b -> mismatch
\node[lbl,left=2mm of t0,yshift=-14mm] {$s=1$};
\node[pcell] at (1*0.65,-1.4) {a};
\node[pcell] at (2*0.65,-1.4) {b};
\node[lbl,right=8mm of t4,yshift=-14mm] {mismatch};
% shift 2: positions 2,3 -> c,a vs a,b -> mismatch
\node[lbl,left=2mm of t0,yshift=-21mm] {$s=2$};
\node[pcell] at (2*0.65,-2.1) {a};
\node[pcell] at (3*0.65,-2.1) {b};
\node[lbl,right=8mm of t4,yshift=-21mm] {mismatch};
% shift 3: positions 3,4 -> a,b -> valid
\node[lbl,left=2mm of t0,yshift=-28mm] {$s=3$};
\node[good] at (3*0.65,-2.8) {a};
\node[good] at (4*0.65,-2.8) {b};
\node[lbl,right=8mm of t4,yshift=-28mm] {valid};
\end{tikzpicture}
\caption{The pattern $P=\texttt{ab}$ slid across $T=\texttt{abcab}$. Each row is
one shift $s$. There are $n-m+1=4$ candidate shifts; the green rows ($s=0$ and
$s=3$) line up exactly, the orange ones do not. String matching is the art of
checking these rows without re-reading the text underneath.}
\label{ch31:fig-shifts}
\end{figure}

\begin{focusbox}
Two numbers govern every analysis in this chapter: $n$, the text length, and
$m$, the pattern length. There are $n-m+1$ candidate shifts. The core challenge
is keeping the cost per shift small, and -- the deeper trick -- avoiding
redundant comparisons by reusing what earlier alignments already revealed.
\end{focusbox}

A word on strings as a tool. We use $\Sigma^*$ for the set of all finite
strings over $\Sigma$, including the empty string $\eps$. If $w$ is a string,
$|w|$ is its length. The \emph{concatenation} of $x$ and $y$, written $xy$, is
$x$ followed by $y$, of length $|x|+|y|$. A string $w$ is a \emph{prefix} of
$x$, written $w \sqsubset x$, if $x = wy$ for some string $y$ -- $w$ is a chunk
off the front. It is a \emph{suffix}, written $w \sqsupset x$, if $x = yw$ for
some $y$ -- a chunk off the back. Thus $\texttt{ab} \sqsubset \texttt{abcab}$
and $\texttt{cab} \sqsupset \texttt{abcab}$. Prefix and suffix are the only two
relations we need for the fastest algorithms in this chapter, so we fix them
firmly now.

\section{The naive matcher: just try every shift}
\label{ch31:naive}

Here is the algorithm you would write without thinking. Slide the pattern to
each of the $n-m+1$ shifts. At each one, compare characters left to right until
either they all match (report the shift) or one of them disagrees (give up,
slide one notch right, start over).

\begin{algorithm}[H]
\caption{\textsc{Naive-String-Matcher}}
\KwIn{Text $T[1\,..\,n]$, pattern $P[1\,..\,m]$.}
\KwOut{Every valid shift $s$.}
\For{$s \leftarrow 0$ \KwTo $n-m$}{
  $j \leftarrow 1$\;
  \While{$j \le m$ \textbf{and} $T[s+j] = P[j]$}{
    $j \leftarrow j + 1$\;
  }
  \If{$j = m+1$}{
    report ``$P$ occurs with shift $s$''\;
  }
}
\end{algorithm}

The inner \textbf{while} marches $j$ forward as long as characters agree. If it
makes it all the way to $j = m+1$, every one of the $m$ characters matched, so
$s$ is valid. Otherwise it bailed early.

\begin{labbox}
Let us trace it. Text $T = \texttt{a b a b a b c a}$ (so $n=8$) and pattern
$P = \texttt{a b a b c}$ (so $m=5$). The candidate shifts run from $0$ to $3$.

\textbf{Shift $s=0$:} compare $T[1]=\texttt{a}$ vs $P[1]=\texttt{a}$ (ok),
$T[2]=\texttt{b}$ vs $P[2]=\texttt{b}$ (ok), $T[3]=\texttt{a}$ vs $P[3]=\texttt{a}$
(ok), $T[4]=\texttt{b}$ vs $P[4]=\texttt{b}$ (ok), $T[5]=\texttt{a}$ vs
$P[5]=\texttt{c}$ -- \emph{mismatch}. Four comparisons matched, the fifth failed.

\textbf{Shift $s=1$:} compare $T[2]=\texttt{b}$ vs $P[1]=\texttt{a}$ --
\emph{mismatch} on the very first character. One comparison.

\textbf{Shift $s=2$:} $T[3]=\texttt{a}$ vs $\texttt{a}$ (ok), $T[4]=\texttt{b}$
vs $\texttt{b}$ (ok), $T[5]=\texttt{a}$ vs $\texttt{a}$ (ok), $T[6]=\texttt{b}$
vs $\texttt{b}$ (ok), $T[7]=\texttt{c}$ vs $\texttt{c}$ (ok). All five match!
Report shift $s=2$.

\textbf{Shift $s=3$:} $T[4]=\texttt{b}$ vs $P[1]=\texttt{a}$ -- mismatch.

The only valid shift is $s=2$, and indeed $T[3\,..\,7]=\texttt{ababc}=P$. Notice
the wasteful part: at shift $0$ we worked through four matching characters before
the fifth failed, and then at shift $1$ we threw all of that away and started
from scratch. That waste is the reason the rest of this chapter exists.
\end{labbox}

\subsection{How slow can the naive matcher be?}

Each shift costs at most $m$ comparisons in the inner loop, and there are
$n-m+1$ shifts, so the total is at most $(n-m+1)\,m$ comparisons. In the worst
case this bound is tight. Take a text of $n$ copies of the letter \texttt{a} and
a pattern of $m$ \texttt{a}'s followed by one \texttt{b}, e.g.
$T=\texttt{aaaaaa}$ and $P=\texttt{aab}$. At every shift the inner loop matches
all the leading \texttt{a}'s and only fails on the final \texttt{b}, so each of
the $n-m+1$ shifts burns the full $m$ comparisons.

\begin{complexitybox}
Naive string matching uses $O((n-m+1)\,m)$ character comparisons, which is
$O(nm)$ in the worst case. It uses $O(1)$ extra space and needs no preprocessing
of the pattern. Its average case on random text is actually decent -- close to
$O(n)$ -- because mismatches usually come fast. But the worst case is genuinely
quadratic, and that is what we are out to beat.
\end{complexitybox}

\begin{trapbox}
The naive matcher's flaw is not that it checks every shift -- in the worst case
every shift really might be valid, so you cannot avoid \emph{looking} at each.
The flaw is that it re-reads text characters it has already examined. At shift
$0$ it read $T[1\,..\,5]$; at shift $1$ it re-read $T[2]$; and so on. A single
text character can be inspected up to $m$ times. The fast algorithms all share
one promise: \emph{look at each text character a bounded number of times} --
ideally just once.
\end{trapbox}

\begin{aestheticbox}
Why care about beating $O(nm)$? Because $n$ is often colossal. A genome is three
billion characters. A web-scale text index runs to terabytes. A log file you are
grepping might be gigabytes. If your pattern is a few dozen characters and your
matcher is quadratic in the bad case, an adversary -- or just unlucky data -- can
turn a one-second search into an hour. Linear-time matching is not a luxury; at
that scale it is the difference between a tool that works and one that hangs.
\end{aestheticbox}

\section{Rabin--Karp: matching by fingerprint}
\label{ch31:rabinkarp}

Here is a different angle, and it starts with a question that sounds unrelated.
Comparing two strings of length $m$ takes up to $m$ character comparisons, but
comparing two \emph{numbers} takes one operation. What if we could turn each
length-$m$ window of text into a single number, turn the pattern into a single
number, and just compare numbers?

That is the Rabin--Karp idea. We assign to every string a numeric
\emph{fingerprint}, computed so that equal strings always get equal
fingerprints. Then, instead of comparing the pattern against each text window
character by character, we compare fingerprints -- one integer comparison per
window.

\subsection{Treating a string as a number}

Suppose for a moment the alphabet is the ten digits $\{0,1,\ldots,9\}$, so
$|\Sigma| = 10$. Then a string \emph{is} a number written in base $10$. The
pattern $P = \texttt{31415}$ has the value
\[
  p = 3\cdot 10^4 + 1\cdot 10^3 + 4\cdot 10^2 + 1\cdot 10^1 + 5\cdot 10^0
    = 31415.
\]
For a general alphabet, treat each character as a ``digit'' in base
$d = |\Sigma|$, and the same construction gives every length-$m$ string a value
between $0$ and $d^m - 1$. Let $t_s$ denote the value of the length-$m$ text
window starting after shift $s$, i.e. of $T[s+1\,..\,s+m]$. Then $s$ is a valid
shift exactly when $t_s = p$ -- provided we computed both values correctly.

There are two problems, and Rabin--Karp solves both with one stroke each.

\textbf{Problem one: the numbers are astronomically large.} A pattern of length
$m=20$ over bytes has a value up to $256^{20}$, a $48$-digit number. Arithmetic
on such numbers is not $O(1)$ -- it is exactly the per-character work we were
trying to escape. \emph{Fix:} do all arithmetic modulo a prime $q$. Choose $q$
small enough that $dq$ fits in one machine word, so every operation is genuinely
$O(1)$. Now $p$ and each $t_s$ are reduced mod $q$ to a small number we call a
\emph{fingerprint}.

\textbf{Problem two: computing each $t_s$ from scratch costs $m$ operations,
which buys us nothing.} \emph{Fix:} compute $t_{s+1}$ from $t_s$ in
\emph{constant} time with a \emph{rolling} update. Slide the window one step
right: you drop the high-order character $T[s+1]$ and append the new low-order
character $T[s+m+1]$. In base $d$ that is
\[
  t_{s+1} = d\,\bigl(t_s - T[s+1]\cdot d^{\,m-1}\bigr) + T[s+m+1].
\]
Read it left to right: subtract off the contribution of the leaving character
($T[s+1]$, which sat in the top place worth $d^{m-1}$), shift everything up one
place by multiplying by $d$, then add the entering character in the ones place.
A car odometer rolls over the same way, and this update is the heart of the
method. Figure~\ref{ch31:fig-roll} shows one such step on the decimal example.

\begin{figure}[H]
\centering
\begin{tikzpicture}[
  font=\small,
  cell/.style={draw=gray!60,fill=gray!5,minimum size=8mm,inner sep=0pt,font=\ttfamily\normalsize},
  lbl/.style={font=\sffamily\footnotesize}
]
% Ruler: single digits spaced generously
\foreach \i in {1,...,7} {
  \node[lbl,text=gray!70] at (\i*1.15 - 1.15, 0.65) {\scriptsize \i};
}
\node[lbl,text=gray!70] at (-0.8, 0.65) {\scriptsize index $i$:};

% Tape cells
\node[cell] (a1) at (0,0) {2};
\node[lbl,left=4mm of a1] {Text $T$:};
\node[cell] (a2) at (1.15,0) {3};
\node[cell] (a3) at (2.3,0) {5};
\node[cell] (a4) at (3.45,0) {9};
\node[cell] (a5) at (4.6,0) {0};
\node[cell] (a6) at (5.75,0) {2};
\node[cell] (a7) at (6.9,0) {3};

% Highlighting windows
\draw[quotientorange,very thick,rounded corners=3pt] (-0.48,-0.48) rectangle (5.08,0.48);
\node[lbl,quotientorange!90!black,align=center] at (2.3, 1.45) {\textbf{Window $s=0$}: $T[1..5] = 23590$\\[1pt]$t_0 \equiv 8 \pmod{13}$};

\draw[targetgreen!80!black,very thick,rounded corners=3pt,densely dashed] (0.67,-0.52) rectangle (6.23,0.42);
\node[lbl,targetgreen!80!black,align=center] at (2.9, -1.8) {\textbf{Window $s=1$}: $T[2..6] = 35902$\\[1pt]$t_1 \equiv 9 \pmod{13}$};

% High digit leaves
\draw[->,>=Stealth,thick,trapred] (a1.south) -- ++(0,-0.45)
  node[below,lbl,trapred,align=center] {drop high digit\\ $-2 \cdot 10^4$};

% Shift and add low digit
\draw[<-,>=Stealth,thick,targetgreen!80!black] (a6.south) -- ++(0,-0.45)
  node[below,lbl,targetgreen!80!black,align=center] {append low digit\\ $+2$};
\end{tikzpicture}
\caption{One rolling step on $T=\texttt{235902\ldots}$ with $m=5$, $d=10$,
$q=13$. The window slides from $T[1\,..\,5]$ (orange) to $T[2\,..\,6]$ (green):
the high digit $T[1]=2$ leaves the top place (worth $10^4$), the remainder shifts
left ($\times 10$), and the new digit $T[6]=2$ enters the units place. The new
fingerprint $t_1=\bigl(10(t_0-T[1]\cdot h)+T[6]\bigr)\bmod 13$ costs $O(1)$ arithmetic
operations rather than $O(m)$ character scans.}
\label{ch31:fig-roll}
\end{figure}

\begin{intuitionbox}
A fingerprint is a tiny summary of a long string -- like the last few digits of
a credit card number, or a checksum. If two strings are equal, their
fingerprints are equal, guaranteed. The converse is \emph{not} guaranteed: two
different strings can collide on the same fingerprint, the way two people can
share a birthday. So a fingerprint match is a \emph{strong hint}, not a proof.
Rabin--Karp uses fingerprints to skip the expensive character-by-character check
almost everywhere, and only does the real check on the rare windows whose
fingerprints actually match the pattern's.
\end{intuitionbox}

\subsection{Modular arithmetic and the spurious-hit check}

Working modulo $q$ means $t_s$ is really $t_s \bmod q$. The valid-shift test
becomes: if $t_s \equiv p \pmod q$, the window is a \emph{candidate}. But
congruence mod $q$ does not imply equality of the underlying strings -- two
different windows can be congruent mod $q$ by coincidence. We call such a window
a \emph{spurious hit}. So whenever the fingerprints match, we must \emph{verify}
by comparing the actual characters $T[s+1\,..\,s+m]$ against $P[1\,..\,m]$. If
they truly match, report $s$; if not, it was a spurious hit and we move on.

\begin{algorithm}[H]
\caption{\textsc{Rabin-Karp-Matcher}}
\KwIn{Text $T[1\,..\,n]$, pattern $P[1\,..\,m]$, radix $d=|\Sigma|$, prime $q$.}
\KwOut{Every valid shift $s$.}
$h \leftarrow d^{\,m-1} \bmod q$ \tcp*{value of the top place}
$p \leftarrow 0$;\quad $t \leftarrow 0$\;
\For{$i \leftarrow 1$ \KwTo $m$}{
  $p \leftarrow (d\cdot p + P[i]) \bmod q$ \tcp*{Horner: fingerprint of $P$}
  $t \leftarrow (d\cdot t + T[i]) \bmod q$ \tcp*{fingerprint of $T[1..m]$}
}
\For{$s \leftarrow 0$ \KwTo $n-m$}{
  \If{$p = t$}{
    \If{$P[1\,..\,m] = T[s+1\,..\,s+m]$}{
      report ``$P$ occurs with shift $s$''\;
    }
  }
  \If{$s < n-m$}{
    $t \leftarrow \bigl(d\,(t - T[s+1]\cdot h) + T[s+m+1]\bigr) \bmod q$\;
  }
}
\end{algorithm}

The first loop is \emph{Horner's rule}: building a number digit by digit, each
step multiplying the running total by the base and adding the next digit, all
mod $q$. The second loop walks the windows, comparing fingerprints, verifying
hits, and rolling the window forward. (In real code, after subtracting
$T[s+1]\cdot h$ the value can go negative; add $q$ before taking the mod to keep
it in range.)

\begin{labbox}
Let us do the classic decimal example. Text
$T = \texttt{2 3 5 9 0 2 3 1 4 1 5 2 6 7 3 9 9 2 1}$ ($n=19$),
pattern $P = \texttt{3 1 4 1 5}$ ($m=5$), radix $d=10$, prime $q=13$.

First the pattern's fingerprint. By Horner mod $13$:
$31415 \bmod 13$. Compute $31415 = 13 \cdot 2416 + 7$, so $p = 7$.

Now the first window $T[1\,..\,5] = \texttt{23590}$:
$23590 = 13\cdot 1814 + 8$, so $t_0 = 8$. Since $8 \ne 7$, no candidate.

Roll forward. The high digit leaving is $T[1]=2$, the place value
$h = 10^4 \bmod 13$. Since $10^4 = 10000 = 13\cdot 769 + 3$, we have $h = 3$.
Each roll: $t' = (10\,(t - T[\text{leave}]\cdot 3) + T[\text{enter}]) \bmod 13$.

The window $\texttt{31415}$ appears as $T[7\,..\,11]$, i.e. at shift $s=6$. Its
true value $31415 \bmod 13 = 7 = p$, so the fingerprints match: a candidate. We
verify $T[7\,..\,11]=\texttt{31415}=P$ -- a real hit. Report $s=6$.

Watch for a spurious hit. The window $\texttt{67399}$ has value
$67399 = 13\cdot 5184 + 7$, so its fingerprint is also $7 = p$! The fingerprints
match, but $\texttt{67399} \ne \texttt{31415}$. Character-by-character
verification catches the spurious collision, so we do \emph{not} report a match. That window was a spurious hit
-- two strings sharing a fingerprint by coincidence, exactly the collision the
intuition box warned about.
\end{labbox}

\subsection{Why it is usually fast, and why it can be slow}

Preprocessing -- computing $p$, $t_0$, and $h$ -- is $\Theta(m)$. The main loop
runs $n-m+1$ times; each iteration does $O(1)$ fingerprint work plus, only on a
hit, an $O(m)$ verification. So the running time hinges entirely on how often
fingerprints match.

In the worst case fingerprints match every time. The cruel input is $T$ all
\texttt{a}'s, $P$ all \texttt{a}'s: every window equals the pattern, every shift
is a real hit, and we pay $\Theta(m)$ to verify each, for $\Theta((n-m+1)m)$ --
right back to naive's quadratic. But that input is degenerate. On ``ordinary''
data, with a well-chosen prime $q$, spurious hits are rare. If we model the
fingerprints as roughly uniform over the $q$ residues, the chance a given
invalid shift looks like a hit is about $1/q$. With at most $n$ shifts, the
expected number of spurious hits is about $n/q$, and the expected verification
work is $O((n/q)\,m)$. If $q$ is chosen larger than $m$ (easy, since $q$ can be
nearly word-sized), this is $O(n)$, and the whole algorithm runs in expected
\[
  O(n) + O(v\cdot m),
\]
where $v$ is the number of valid shifts -- linear, plus the unavoidable cost of
confirming the real occurrences.

\begin{complexitybox}
\textbf{Rabin--Karp.} Preprocessing $\Theta(m)$. Worst-case matching
$\Theta((n-m+1)m)$ -- no better than naive, attained when fingerprints collide
constantly. Expected matching $O(n + vm)$ on data where collisions are rare and
$q$ is large, which for a small fixed number of occurrences is $O(n)$. Extra
space $O(1)$.
\end{complexitybox}

\begin{focusbox}
Beyond single-pattern search, fingerprinting scales naturally: because a fingerprint is
just an integer, you can keep a \emph{set} of target fingerprints in a hash table
and search for many patterns simultaneously in one pass, or extend the rolling-hash
to two dimensions to locate subimages inside large bitmaps. The same rolling hash
powers synchronization utilities like \texttt{rsync} and content-defined storage deduplication.
CLRS covers Rabin--Karp in Section~32.2.
\end{focusbox}

\begin{walkbox}
Rabin--Karp gets its speed from a probabilistic gamble: it bets that
fingerprints rarely collide. But an adversary who \emph{knows your prime $q$} can
craft a text where every window collides, dragging you back to $O(nm)$. How would
you defend against such an adversary? (Hint: what if $q$, or the radix, were
chosen at random and kept secret -- how does that change the analysis, and what
does it cost you?) This is the same randomization idea that rescued quicksort and
hashing.
\end{walkbox}

\section{String matching with finite automata}
\label{ch31:automata}

The next idea drops probability entirely and returns to a worst-case guarantee.
It also achieves the goal set out in the trap box above: \emph{read each text
character exactly once, never backing up.} The price is some work up front -- we
compile the pattern into a little machine.

\subsection{What a string-matching automaton is}

A \emph{finite automaton} is a tiny processor with a finite set of \emph{states}.
It starts in a designated start state, then reads the text one character at a
time, and each character makes it jump to a new state according to a fixed
\emph{transition function} $\delta$. Some states are marked \emph{accepting}.
For string matching we build the machine so that, after reading $T[1\,..\,i]$,
its current state is precisely \emph{how much of the pattern we have matched so
far} -- a number from $0$ to $m$. State $m$ is accepting: reaching it means the
last $m$ characters we read were exactly the pattern, so we have just found an
occurrence ending at position $i$.

\begin{definition}[String-matching automaton]
\label{ch31:def-automaton}
For a pattern $P[1\,..\,m]$ the matching automaton has state set
$Q = \{0, 1, \ldots, m\}$, start state $0$, single accepting state $m$, and
transition function $\delta(q, a)$ defined for a state $q$ and a character
$a \in \Sigma$ by
\[
  \delta(q, a) = \text{the length of the longest prefix of } P
  \text{ that is also a suffix of } P_q\,a,
\]
where $P_q = P[1\,..\,q]$ is the length-$q$ prefix of the pattern and $P_q\,a$
means $P_q$ with the character $a$ appended.
\end{definition}

That definition packs in a lot, so let us take it apart. The state
$q$ means ``the last $q$ characters I read were $P[1\,..\,q]$.'' Now a new
character $a$ arrives. The string I have matched, with $a$ stuck on the end, is
$P_q\,a$. The new state must answer: \emph{after seeing $a$, what is the longest
prefix of $P$ I can still claim to have just matched at the end?} That is the
longest prefix of $P$ that is a suffix of $P_q\,a$. If $a$ happened to be exactly
$P[q+1]$, then $P_q\,a = P[1\,..\,q+1]$ is itself a prefix of $P$, and we advance
to state $q+1$. If $a$ was wrong, we cannot advance -- but we might not have to
fall all the way to $0$, because a shorter prefix of $P$ might still match at the
tail. The automaton precomputes the correct landing state for every $(q, a)$ pair.

\begin{intuitionbox}
Think of climbing a ladder with $m$ rungs, where rung $q$ means ``I have matched
the first $q$ characters of the pattern.'' Each text character you read either
moves you up one rung (you read the character the pattern wanted next) or knocks
you down to some lower rung (you read something else). The interesting part is
exactly \emph{how far down} a wrong character knocks you. It does not necessarily
send you back to the ground, because the characters you just read might form the
start of a fresh match. The transition function $\delta$ has the right landing
rung stored for every situation, so the machine never hesitates.
\end{intuitionbox}

Once $\delta$ is built, scanning could hardly be simpler.

\begin{algorithm}[H]
\caption{\textsc{Finite-Automaton-Matcher}}
\KwIn{Text $T[1\,..\,n]$, transition function $\delta$, pattern length $m$.}
\KwOut{Every valid shift $s$.}
$q \leftarrow 0$ \tcp*{start state: nothing matched yet}
\For{$i \leftarrow 1$ \KwTo $n$}{
  $q \leftarrow \delta(q, T[i])$ \tcp*{one table lookup, no backtracking}
  \If{$q = m$}{
    report ``$P$ occurs with shift $s = i - m$''\;
  }
}
\end{algorithm}

The whole body is one lookup per text character and a check for the accepting
state. There is no inner loop, and the pointer $i$ never moves backward. This is
the structure of \texttt{grep} on a fixed string, and of the lexical analyzer in
every compiler.

\subsection{A worked machine}

Let $P = \texttt{ababaca}$, so $m = 7$, over the alphabet
$\Sigma = \{\texttt{a}, \texttt{b}, \texttt{c}\}$. The states are $0$ through
$7$. We compute $\delta(q,a)$ for each by the definition: the longest prefix of
$P$ that is a suffix of $P_q a$.

Take $\delta(5, \texttt{b})$. State $5$ means we have matched
$P_5 = \texttt{ababa}$. Append \texttt{b}: we get $\texttt{ababab}$. What is the
longest prefix of \texttt{ababaca} that is a suffix of \texttt{ababab}? Check
suffixes longest-first: \texttt{ababab}? not a prefix of $P$.
\texttt{babab}? no. \texttt{abab}? yes -- \texttt{abab} is a prefix of $P$
(it is $P_4$). So $\delta(5, \texttt{b}) = 4$. We read a wrong character and only
slipped from rung $5$ to rung $4$, not to the ground.

Take $\delta(5, \texttt{a})$. Append \texttt{a} to \texttt{ababa}: get
\texttt{ababaa}. Longest prefix of $P$ that is a suffix? \texttt{a} works
($P_1$); \texttt{aa} is not a prefix of $P$. So $\delta(5, \texttt{a}) = 1$.

And the happy case $\delta(5, \texttt{c})$. Append \texttt{c}: \texttt{ababac},
which is $P_6$ itself -- a prefix of $P$. So $\delta(5, \texttt{c}) = 6$: we
advance. Here is the full table:

\begin{center}
\begin{tabular}{c|ccc}
\toprule
state $q$ & $\delta(q,\texttt{a})$ & $\delta(q,\texttt{b})$ & $\delta(q,\texttt{c})$ \\
\midrule
$0$ & $1$ & $0$ & $0$ \\
$1$ & $1$ & $2$ & $0$ \\
$2$ & $3$ & $0$ & $0$ \\
$3$ & $1$ & $4$ & $0$ \\
$4$ & $5$ & $0$ & $0$ \\
$5$ & $1$ & $4$ & $6$ \\
$6$ & $7$ & $0$ & $0$ \\
$7$ & $1$ & $2$ & $0$ \\
\bottomrule
\end{tabular}
\end{center}

Here is the same machine as a state diagram. The forward spine (solid blue)
advances character by character along the pattern $P=\texttt{ababaca}$. The
curved fallback transitions (dashed orange) show the non-trivial retreats when
a mismatch occurs, illustrating how the machine preserves the longest matching
prefix without ever rescanning earlier text characters.

\begin{center}
\begin{tikzpicture}[
  >=Stealth,
  state/.style={circle,draw=sourceblue!80!black,fill=sourceblue!10,thick,minimum size=8mm,inner sep=1pt,font=\small},
  accept/.style={circle,draw=targetgreen!80!black,fill=targetgreen!15,very thick,minimum size=8mm,double,inner sep=1pt,font=\small},
  fwd/.style={draw=sourceblue!80!black,->,very thick},
  back/.style={draw=quotientorange!80!black,->,thick,dashed}
]
\node[state] (q0) at (0,0) {$0$};
\node[state] (q1) at (1.5,0) {$1$};
\node[state] (q2) at (3.0,0) {$2$};
\node[state] (q3) at (4.5,0) {$3$};
\node[state] (q4) at (6.0,0) {$4$};
\node[state] (q5) at (7.5,0) {$5$};
\node[state] (q6) at (9.0,0) {$6$};
\node[accept] (q7) at (10.5,0) {$7$};

% Start arrow
\draw[->,thick] (-0.8,0) -- (q0);

% Forward spine
\path
 (q0) edge[fwd] node[above,font=\scriptsize\bfseries]{a} (q1)
 (q1) edge[fwd] node[above,font=\scriptsize\bfseries]{b} (q2)
 (q2) edge[fwd] node[above,font=\scriptsize\bfseries]{a} (q3)
 (q3) edge[fwd] node[above,font=\scriptsize\bfseries]{b} (q4)
 (q4) edge[fwd] node[above,font=\scriptsize\bfseries]{a} (q5)
 (q5) edge[fwd] node[above,font=\scriptsize\bfseries]{c} (q6)
 (q6) edge[fwd] node[above,font=\scriptsize\bfseries]{a} (q7);

% Key fallback transitions
\draw[back] (q1) to[out=120,in=60,looseness=4] node[above,font=\scriptsize]{a} (q1);
\draw[back] (q5) to[bend left=38] node[below,font=\scriptsize]{b} (q4);
\draw[back] (q5) to[bend left=48] node[below,font=\scriptsize]{a} (q1);
\draw[back] (q7) to[bend left=55] node[below,font=\scriptsize]{b} (q2);
\draw[back] (q7) to[bend left=62] node[below,font=\scriptsize]{a} (q1);
\end{tikzpicture}
\end{center}

\begin{labbox}
Let us run the machine on $T = \texttt{a b a b a b a c a b a}$ ($n=11$). Start in
state $q=0$. Read each character and look up $\delta$.

\begin{center}
\begin{tabular}{c c c c}
\toprule
$i$ & $T[i]$ & new state $q=\delta(q,T[i])$ & note \\
\midrule
$1$ & a & $\delta(0,a)=1$ & \\
$2$ & b & $\delta(1,b)=2$ & \\
$3$ & a & $\delta(2,a)=3$ & \\
$4$ & b & $\delta(3,b)=4$ & \\
$5$ & a & $\delta(4,a)=5$ & \\
$6$ & b & $\delta(5,b)=4$ & wrong char, slip to $4$ \\
$7$ & a & $\delta(4,a)=5$ & \\
$8$ & c & $\delta(5,c)=6$ & \\
$9$ & a & $\delta(6,a)=7$ & \textbf{accept!} shift $s=9-7=2$ \\
$10$ & b & $\delta(7,b)=2$ & \\
$11$ & a & $\delta(2,a)=3$ & \\
\bottomrule
\end{tabular}
\end{center}

At $i=9$ the machine entered state $7=m$, so the pattern $\texttt{ababaca}$ occurs
with shift $2$ -- check: $T[3\,..\,9]=\texttt{ababaca}=P$. Notice especially row
$i=6$: a stray \texttt{b} where the pattern wanted \texttt{c}. The naive matcher
would have restarted; the automaton simply slipped from state $5$ to state $4$ and
kept reading $T$ forward, never re-examining a single character.
\end{labbox}

\subsection{Cost, and the one catch}

Scanning is cheap: the loop runs $n$ times with one table lookup each, so
matching is $\Theta(n)$ -- guaranteed, no probability, no backtracking. The catch
is building $\delta$. The table has $m+1$ states times $|\Sigma|$ characters
$= (m+1)|\Sigma|$ entries, and the straightforward way to compute each entry tests
prefixes against suffixes, costing up to $O(m)$ per entry, for $O(m^2|\Sigma|)$
preprocessing. (A cleverer construction, closely related to the prefix function
of the next section, brings this down to $O(m|\Sigma|)$.)

\begin{complexitybox}
\textbf{Finite-automaton matcher.} Matching $\Theta(n)$ in the worst case as
well as on average. Preprocessing $O(m^2|\Sigma|)$ by the direct
construction, $O(m|\Sigma|)$ by the more structured one. Extra space $O(m|\Sigma|)$ for
the transition table. The matching phase reads each text character exactly once
and never backtracks.
\end{complexitybox}

\begin{trapbox}
The automaton's transition table is $\Theta(m|\Sigma|)$ in size. For DNA
($|\Sigma|=4$) or ASCII this is compact. But for Unicode, where $|\Sigma|$ can
be over a million, the table becomes enormous -- it can dwarf the text you are
searching. This dependence on $|\Sigma|$ is the automaton's primary drawback, and it
is exactly the weakness KMP, our final algorithm, was designed to remove.
\end{trapbox}

\section{The Knuth--Morris--Pratt algorithm}
\label{ch31:kmp}

KMP achieves the same $\Theta(n)$ scan as the automaton, with the same
no-backtracking guarantee, but its preprocessing is $\Theta(m)$ and
\emph{entirely independent of the alphabet size}. It needs no transition table at all.
Instead it carries one small array, the \emph{prefix function}, which captures
the pattern's internal self-similarity. You can think of the prefix function as
the automaton's table compressed down to what actually matters.

\subsection{The insight: a mismatch already tells you a lot}

Return once more to the wasteful moment in the naive matcher. We were matching
$P$ against the text, we had matched $q$ characters, and then character $q+1$
disagreed. The naive matcher slides the pattern one step and restarts from the
beginning. But consider what we already \emph{know}: the last $q$ text characters
were exactly $P[1\,..\,q]$. That is real information about the text, paid for with
$q$ comparisons, and KMP refuses to throw it away.

Here is the question KMP asks. We matched $P[1\,..\,q]$ and then failed. We must
slide the pattern forward -- but how far can we slide it without missing a
possible match? A match cannot start anywhere that would contradict the $q$
characters we already verified. The pattern, after sliding, must agree with the
tail of $P[1\,..\,q]$ wherever they now overlap. In other words: we want the
longest \emph{proper} prefix of $P[1\,..\,q]$ that is also a suffix of
$P[1\,..\,q]$. That length tells us the new state to drop to -- \emph{and it
depends only on the pattern, not on the text at all.}

\begin{definition}[Prefix function]
\label{ch31:def-prefix}
The \emph{prefix function} of a pattern $P[1\,..\,m]$ is the array
$\pi[1\,..\,m]$ where
\[
  \pi[q] = \text{the length of the longest proper prefix of } P[1\,..\,q]
  \text{ that is also a suffix of } P[1\,..\,q].
\]
``Proper'' means strictly shorter than $P[1\,..\,q]$ itself (so we never pick the
whole thing). When no proper prefix is also a suffix, $\pi[q] = 0$.
\end{definition}

\begin{intuitionbox}
The prefix function measures how much the pattern \emph{repeats its own
beginning} inside itself. Suppose you have matched the first $q$ characters and
then hit a wall. Those $q$ characters are $P[1\,..\,q]$. If the \emph{end} of
that chunk happens to repeat the \emph{beginning} of the pattern -- say the
last $3$ characters of what you matched are also the first $3$ characters of $P$
-- then you have effectively already matched $3$ characters of a \emph{fresh}
attempt, for free, without re-reading any text. So you do not restart from zero;
you restart from $3$. The prefix function precomputes, for every $q$, exactly how
big that free head start is.
\end{intuitionbox}

\subsection{Computing the prefix function by hand}

Let $P = \texttt{ababaca}$ again, $m=7$. We fill $\pi[q]$ for $q=1,\ldots,7$.

\begin{itemize}
\item $\pi[1]$: $P[1\,..\,1]=\texttt{a}$. A proper prefix is strictly shorter, so
the only candidate is the empty string. $\pi[1]=0$.
\item $\pi[2]$: $P[1\,..\,2]=\texttt{ab}$. Proper prefixes: \texttt{a}. Is
\texttt{a} a suffix of \texttt{ab}? No (the suffix of length $1$ is \texttt{b}).
$\pi[2]=0$.
\item $\pi[3]$: $\texttt{aba}$. Proper prefixes \texttt{ab}, \texttt{a}. Is
\texttt{ab} a suffix? The length-$2$ suffix is \texttt{ba}; no. Is \texttt{a} a
suffix? The length-$1$ suffix is \texttt{a}; yes. $\pi[3]=1$.
\item $\pi[4]$: $\texttt{abab}$. Is \texttt{aba} a suffix? Length-$3$ suffix is
\texttt{bab}; no. Is \texttt{ab} a suffix? Length-$2$ suffix is \texttt{ab};
yes. $\pi[4]=2$.
\item $\pi[5]$: $\texttt{ababa}$. Is \texttt{abab} a suffix? Length-$4$ suffix is
\texttt{baba}; no. Is \texttt{aba} a suffix? Length-$3$ suffix \texttt{aba};
yes. $\pi[5]=3$.
\item $\pi[6]$: $\texttt{ababac}$. Length-$5$ suffix \texttt{babac}? no.
Length-$4$ \texttt{abac}? no. Length-$3$ \texttt{bac}? no. Length-$2$
\texttt{ac}? no. Length-$1$ \texttt{c}? no. $\pi[6]=0$.
\item $\pi[7]$: $\texttt{ababaca}$. Down the suffixes\ldots only \texttt{a}
(length $1$) is a prefix and a suffix. $\pi[7]=1$.
\end{itemize}

\begin{center}
\begin{tabular}{c|ccccccc}
\toprule
$q$ & $1$ & $2$ & $3$ & $4$ & $5$ & $6$ & $7$ \\
$P[q]$ & a & b & a & b & a & c & a \\
$\pi[q]$ & $0$ & $0$ & $1$ & $2$ & $3$ & $0$ & $1$ \\
\bottomrule
\end{tabular}
\end{center}

Read $\pi[5]=3$ aloud: ``if I have matched \texttt{ababa} and then fail, the
last $3$ characters \texttt{aba} are also the first $3$, so I restart from a head
start of $3$.'' That single number is what the automaton needed a whole table
column to express. Figure~\ref{ch31:fig-prefix} draws this overlap: the longest
proper prefix of \texttt{ababa} that is also a suffix.

\begin{figure}[H]
\centering
\begin{tikzpicture}[
  font=\ttfamily\small,
  cell/.style={draw,minimum size=6.5mm,inner sep=0pt},
  base/.style={cell,fill=gray!8},
  pre/.style={cell,fill=sourceblue!22},
  suf/.style={cell,fill=targetgreen!28},
  both/.style={cell,fill=structurepurple!28},
  lbl/.style={font=\sffamily\scriptsize}
]
% P[1..5] = a b a b a
\foreach \c [count=\i from 1] in {a,b,a,b,a} {
  \node[base] (p\i) at (\i*0.75,0) {\c};
}
\node[lbl,left=2mm of p1] {$P[1\,..\,5]$};
% prefix aba = positions 1,2,3 (blue brace above)
\draw[decorate,decoration={brace,amplitude=4pt},sourceblue,thick]
  ($(p1.north west)+(0,3pt)$) -- ($(p3.north east)+(0,3pt)$)
  node[midway,above=4pt,lbl,sourceblue]{prefix \texttt{aba}};
% suffix aba = positions 3,4,5 (green brace below)
\draw[decorate,decoration={brace,amplitude=4pt,mirror},targetgreen,thick]
  ($(p3.south west)+(0,-3pt)$) -- ($(p5.south east)+(0,-3pt)$)
  node[midway,below=4pt,lbl,targetgreen]{suffix \texttt{aba}};
\end{tikzpicture}
\caption{Why $\pi[5]=3$ for $P=\texttt{ababaca}$. Inside $P[1\,..\,5]=\texttt{ababa}$,
the prefix \texttt{aba} (blue, positions $1$--$3$) equals the suffix \texttt{aba}
(green, positions $3$--$5$); they even overlap at position $3$. That length-$3$
agreement is the free head start KMP keeps after a mismatch at state $5$.}
\label{ch31:fig-prefix}
\end{figure}

\begin{figure}[H]
\centering
\begin{tikzpicture}[
  font=\small,
  >=Stealth,
  state/.style={circle,draw=sourceblue!80!black,fill=sourceblue!10,thick,minimum size=8mm,inner sep=0pt},
  fwd/.style={draw=sourceblue!80!black,->,very thick},
  fail/.style={draw=trapred!80!black,->,thick,dashed},
  lbl/.style={font=\sffamily\scriptsize}
]
% States 0 to 7
\foreach \s in {0,...,7} {
  \node[state] (s\s) at (\s*1.55, 0) {$\s$};
}

% Forward matching characters
\path
  (s0) edge[fwd] node[above,lbl,font=\ttfamily\bfseries]{a} (s1)
  (s1) edge[fwd] node[above,lbl,font=\ttfamily\bfseries]{b} (s2)
  (s2) edge[fwd] node[above,lbl,font=\ttfamily\bfseries]{a} (s3)
  (s3) edge[fwd] node[above,lbl,font=\ttfamily\bfseries]{b} (s4)
  (s4) edge[fwd] node[above,lbl,font=\ttfamily\bfseries]{a} (s5)
  (s5) edge[fwd] node[above,lbl,font=\ttfamily\bfseries]{c} (s6)
  (s6) edge[fwd] node[above,lbl,font=\ttfamily\bfseries]{a} (s7);

% Failure links pi[q]
\draw[fail] (s1) to[bend left=32] node[below,lbl,trapred]{$\pi=0$} (s0);
\draw[fail] (s2) to[bend left=38] node[below,lbl,trapred]{$\pi=0$} (s0);
\draw[fail] (s3) to[bend left=32] node[below,lbl,trapred]{$\pi=1$} (s1);
\draw[fail] (s4) to[bend left=32] node[below,lbl,trapred]{$\pi=2$} (s2);
\draw[fail] (s5) to[bend left=32] node[below,lbl,trapred]{$\pi=3$} (s3);
\draw[fail] (s6) to[bend left=42] node[below,lbl,trapred]{$\pi=0$} (s0);
\draw[fail] (s7) to[bend left=45] node[below,lbl,trapred]{$\pi=1$} (s1);

% Annotations above and below
\node[lbl,sourceblue!80!black,above=8pt of s3] {solid blue: advance on match $P[q+1]$};
\node[lbl,trapred,below=22pt of s4] {dashed red: fallback $\pi[q]$ on mismatch};
\end{tikzpicture}
\caption{The KMP prefix function $\pi$ viewed as state failure links. When matching
at state $q=5$ ($P[1\,..\,5]=\texttt{ababa}$) fails to match the next character
\texttt{c}, the algorithm retreats along the dashed link to $\pi[5]=3$, where
the prefix \texttt{aba} is guaranteed to match the text so far. The text pointer
$i$ never steps backward.}
\label{ch31:fig-pi-links}
\end{figure}

\subsection{Computing the prefix function cleverly}

Filling $\pi$ by brute force, as above, costs $O(m^2)$. There is a beautiful
linear way, and it uses the prefix function to compute \emph{itself}. Process
$q$ from $2$ upward, carrying a variable $k = \pi[q-1]$, the length matched so
far. To extend to position $q$, we want to match $P[q]$ against $P[k+1]$. If they
agree, the matched length grows: $k \leftarrow k+1$. If they disagree, we fall
back to the next-shorter candidate, which is precisely $\pi[k]$ -- the prefix
function calling on itself -- and try again, until either we match or $k$ hits
$0$.

\begin{algorithm}[H]
\caption{\textsc{Compute-Prefix-Function}}
\KwIn{Pattern $P[1\,..\,m]$.}
\KwOut{The prefix-function array $\pi[1\,..\,m]$.}
$\pi[1] \leftarrow 0$\;
$k \leftarrow 0$ \tcp*{length of the current longest prefix-suffix}
\For{$q \leftarrow 2$ \KwTo $m$}{
  \While{$k > 0$ \textbf{and} $P[k+1] \ne P[q]$}{
    $k \leftarrow \pi[k]$ \tcp*{fall back to the next-best candidate}
  }
  \If{$P[k+1] = P[q]$}{
    $k \leftarrow k + 1$ \tcp*{extend the match by one}
  }
  $\pi[q] \leftarrow k$\;
}
\Return $\pi$\;
\end{algorithm}

\subsection{KMP itself}

With $\pi$ in hand, matching is the automaton's spirit without the automaton's
table. We scan the text with a pointer $i$ that never retreats, and a state $q$
counting how much of the pattern is currently matched. On a mismatch we do not
restart -- we consult $\pi[q]$ to find the right fallback state, exactly as
\textsc{Compute-Prefix-Function} fell back on itself.

\begin{algorithm}[H]
\caption{\textsc{KMP-Matcher}}
\KwIn{Text $T[1\,..\,n]$, pattern $P[1\,..\,m]$.}
\KwOut{Every valid shift $s$.}
$\pi \leftarrow \textsc{Compute-Prefix-Function}(P)$\;
$q \leftarrow 0$ \tcp*{number of pattern characters matched so far}
\For{$i \leftarrow 1$ \KwTo $n$}{
  \While{$q > 0$ \textbf{and} $P[q+1] \ne T[i]$}{
    $q \leftarrow \pi[q]$ \tcp*{mismatch: slide using the prefix function}
  }
  \If{$P[q+1] = T[i]$}{
    $q \leftarrow q + 1$ \tcp*{match: advance the state}
  }
  \If{$q = m$}{
    report ``$P$ occurs with shift $s = i - m$''\;
    $q \leftarrow \pi[q]$ \tcp*{look for the next, possibly overlapping, match}
  }
}
\end{algorithm}

The two loops are siblings. \textsc{Compute-Prefix-Function} matches the pattern
against \emph{itself}; \textsc{KMP-Matcher} matches the pattern against the
\emph{text}. Both advance on agreement and both consult $\pi$ on disagreement.
After a full match we set $q \leftarrow \pi[m]$ so that overlapping occurrences
(think of pattern \texttt{aa} in text \texttt{aaa}) are not missed.

\begin{labbox}
Run KMP with $P = \texttt{ababaca}$ (and $\pi = [0,0,1,2,3,0,1]$ from above) on
$T = \texttt{a b a b a b a c a b a}$ ($n=11$). Start $q=0$.

\begin{center}
\begin{tabular}{c c l c}
\toprule
$i$ & $T[i]$ & action & $q$ after \\
\midrule
$1$ & a & $P[1]=$a matches; advance & $1$ \\
$2$ & b & $P[2]=$b matches; advance & $2$ \\
$3$ & a & $P[3]=$a matches; advance & $3$ \\
$4$ & b & $P[4]=$b matches; advance & $4$ \\
$5$ & a & $P[5]=$a matches; advance & $5$ \\
$6$ & b & $P[6]=$c $\ne$ b: fall back $q\!\leftarrow\!\pi[5]=3$; now $P[4]=$b matches; advance & $4$ \\
$7$ & a & $P[5]=$a matches; advance & $5$ \\
$8$ & c & $P[6]=$c matches; advance & $6$ \\
$9$ & a & $P[7]=$a matches; advance $\to 7=m$: \textbf{report }$s=2$; set $q\!\leftarrow\!\pi[7]=1$ & $1$ \\
$10$ & b & $P[2]=$b matches; advance & $2$ \\
$11$ & a & $P[3]=$a matches; advance & $3$ \\
\bottomrule
\end{tabular}
\end{center}

Study row $i=6$. We had matched $5$ characters (\texttt{ababa}) and the text gave
us \texttt{b} where the pattern wanted \texttt{c}. Instead of crawling back, KMP
consulted $\pi[5]=3$: the prefix \texttt{aba} of what we matched is also its
suffix, so we keep a head start of $3$ and immediately try $P[4]=\texttt{b}$
against the same $T[6]=\texttt{b}$ -- it matches, and we are at state $4$ without
having moved $i$ backward at all. That single row contains the entire idea.
\end{labbox}

\subsection{Why KMP is linear: an amortized argument}

The worry with KMP is the inner \textbf{while} loop. Could it spin many times per
character and blow up the cost? The answer is a clean amortized argument that
rests on watching the state $q$ like a bank balance.

\begin{focusbox}
Watch the value of $q$ across the whole run. Each iteration of the outer
\textbf{for} loop increases $q$ by \emph{at most $1$} (the single ``advance''
line). Each iteration of the inner \textbf{while} loop \emph{strictly decreases}
$q$ (since $\pi[q] < q$ always -- the prefix is \emph{proper}). And $q$ never
goes below $0$. So $q$ is a counter that rises by at most $1$ per text character
and can never fall below zero.
\end{focusbox}

\begin{formalbox}
Use the \emph{potential method}. Let the potential be the current value of $q$,
which starts at $0$ and is always $\ge 0$. Across the $n$ iterations of the outer
loop, the total \emph{increase} in $q$ is at most $n$ (at most $1$ per
iteration). Every inner-\textbf{while} iteration decreases $q$ by at least $1$.
Since total decrease cannot exceed total increase plus the starting value, the
inner loop executes at most $n$ times \emph{over the entire algorithm} -- not
per character, but in total. Hence the matching phase does $O(n)$ work. Computing
$\pi$ is the same argument over $P$, giving $O(m)$. Together, KMP runs in
$\Theta(n + m) = \Theta(n)$ time.
\end{formalbox}

This is the same accounting trick that tames the dynamic array and the union-find
structure: a single operation may be expensive, but expensive steps must be
``paid for'' by cheap ones that came before, so the average over the whole run is
small. Here, every fall-back must have been earned by an earlier advance.

Figure~\ref{ch31:fig-potential} plots $q$ across the trace from the lab above.
Every up-step is one advance; every down-step is a fall-back. Because $q$ can
only fall as far as it has previously climbed, the total descent over the whole
run is bounded by the total ascent, which is at most $n$.

\begin{figure}[H]
\centering
\begin{tikzpicture}[font=\sffamily\scriptsize,x=8.5mm,y=5mm]
% axes
\draw[->] (0.5,0) -- (12,0) node[right]{$i$};
\draw[->] (0.5,0) -- (0.5,7.3) node[above]{$q$};
\foreach \y in {1,...,7} \draw (0.4,\y) -- (0.5,\y) node[left,inner sep=1pt]{\y};
\foreach \x in {1,...,11} \node[below] at (\x,0) {\x};
% q values after each i: 1,2,3,4,5,4,5,6,1,2,3
\def\qs{1/1,2/2,3/3,4/4,5/5,6/4,7/5,8/6,9/1,10/2,11/3}
% draw step line from (0,0)
\draw[sourceblue,very thick]
  (0.5,0) -- (1,1) -- (2,2) -- (3,3) -- (4,4) -- (5,5)
  -- (5,4) -- (6,4)          % i=6 fall-back 5->4 then advance to 4
  -- (7,5) -- (8,6)
  -- (9,7)                   % i=9 advance to 7 (match)
  -- (9,1)                   % reset q<-pi[7]=1
  -- (10,2) -- (11,3);
% mark the match
\fill[targetgreen] (9,7) circle (2pt);
\node[targetgreen,above right] at (9,7) {match, $q=m$};
% mark a fall-back
\draw[<-, trapred, thick] (5.0, 4.6) -- ++(-0.6, 0.6) node[above, font=\sffamily\scriptsize, trapred] {fall-back};
\end{tikzpicture}
\caption{The state $q$ during the KMP run on $T=\texttt{abababacaba}$ with
$P=\texttt{ababaca}$. Diagonal rises are advances (at most one per text
character); vertical drops are fall-backs through $\pi$ (at $i=6$, and the reset
to $\pi[7]=1$ after the match at $i=9$). Total drop $\le$ total rise $\le n$, so
the inner loop fires $O(n)$ times in all.}
\label{ch31:fig-potential}
\end{figure}

\begin{complexitybox}
\textbf{Knuth--Morris--Pratt.} Preprocessing $\Theta(m)$, and -- unlike the
automaton -- \emph{independent of $|\Sigma|$}. Matching $\Theta(n)$ in the worst
case. Total $\Theta(n+m)$. Extra space $\Theta(m)$ for the array $\pi$. Each text
character is examined a small \emph{amortized} constant number of times; the
pointer $i$ never moves backward.
\end{complexitybox}

\begin{focusbox}
The elegance of Knuth--Morris--Pratt lies in extracting only what a full automaton
truly requires: rather than storing transitions for every character in the alphabet,
the prefix function $\pi$ records the single best fallback state on mismatch.
This \emph{failure function} principle is foundational across parsing and automata theory,
directly inspiring the Aho--Corasick algorithm for matching dictionaries of multiple patterns simultaneously.
CLRS treats finite-automaton matching in Section~32.3 and KMP in Section~32.4;
the prefix function there is named exactly as we have named it here ($\pi$).
\end{focusbox}

\begin{trapbox}
A common confusion: people think the inner \textbf{while} loop makes KMP
quadratic because ``it has a loop inside a loop.'' It does not. The number of
nested-loop iterations is bounded \emph{globally}, not per outer step, by the
potential argument. Counting loops by their nesting depth, rather than by how
many times they truly run, is one of the most frequent analysis errors in all of
algorithms -- and KMP is the classic cautionary tale.
\end{trapbox}

\section{Putting the four side by side}
\label{ch31:summary}

We have gone from the obvious method to an optimal one. Here is the comparison
side by side.

\begin{center}
\begin{tabular}{l l l l}
\toprule
Algorithm & Preprocessing & Matching & Notes \\
\midrule
Naive & $0$ & $O((n-m+1)m)$ & no extra space; re-reads text \\
Rabin--Karp & $\Theta(m)$ & $O(n+vm)$ expected & worst case $O((n-m+1)m)$ \\
Finite automaton & $O(m|\Sigma|)$ & $\Theta(n)$ & big table; $|\Sigma|$-dependent \\
KMP & $\Theta(m)$ & $\Theta(n)$ & small array; $|\Sigma|$-independent \\
\bottomrule
\end{tabular}
\end{center}

Here $v$ is the number of valid shifts. Notice the arc of the chapter. Naive is
simplest but re-reads text. Rabin--Karp escapes re-reading on \emph{average} by
gambling on fingerprints. The automaton escapes it \emph{always} but pays in a
fat, alphabet-sized table. KMP keeps the automaton's iron-clad linear guarantee
while shrinking the table down to a single $m$-element array that does not care
how big the alphabet is. That progression -- obvious, then probabilistic, then
guaranteed-but-heavy, then guaranteed-and-light -- recurs all over algorithm
design.

\begin{aestheticbox}
Exact string matchers run inside everyday tools: KMP and Boyer--Moore variants
power \texttt{grep}, your editor's find, DNA aligners scanning genomes,
inside network intrusion detectors comparing packets against attack signatures at
line rate, and inside the tokenizers of compilers. The automaton view is the
backbone of regular-expression engines. And Rabin--Karp's rolling hash is the
engine of file-syncing and deduplication tools that decide which blocks of a
multi-gigabyte file actually changed. The ideas in this short chapter run quietly
underneath a large fraction of everyday text processing.
\end{aestheticbox}

\subsection{From text streams to genomic search}

The elimination of backtracking did not just save CPU cycles; it transformed what
was computationally possible:
\begin{itemize}
  \item \textbf{True streaming search without buffering.} Because KMP inspects
        each character of the text exactly once and never rewinds the text
        pointer, it can operate over infinite, non-seekable data streams. Network
        intrusion detection systems (like Snort) inspect packet payloads at line
        rate without caching gigabytes of traffic in RAM.
  \item \textbf{Sublinear skipping in practice.} The Boyer--Moore algorithm (1977)
        showed that scanning the pattern from right to left allows skipping vast
        swaths of text without reading them at all. On English prose, it
        routinely achieves an effective running time of $O(n/m)$ character
        reads, making it the practical engine under Unix \texttt{grep}.
  \item \textbf{Multi-pattern scaling.} Generalizing the finite-automaton
        view of Section~\ref{ch31:automata} to a tree of patterns yields the
        \emph{Aho--Corasick algorithm} (1975). It matches an entire dictionary
        of $k$ patterns simultaneously in a single linear pass over the text,
        with running time $O(n + m_{\text{total}} + z)$ where $z$ is the number
        of matches.
  \item \textbf{Algebraic fingerprinting and chunking.} Rabin--Karp's rolling
        hash broke out of string matching entirely. It became the backbone of
        content-defined chunking (e.g., the \texttt{rsync} protocol and backup
        deduplication systems), where files are divided into variable-length
        blocks at fingerprint boundaries to synchronize differences across the
        internet with minimal bandwidth.
\end{itemize}

\subsection{Practical trade-offs and engineering limits}

Despite their mathematical elegance, these algorithms have sharp boundaries where
they encounter practical constraints:
\begin{itemize}
  \item \textbf{Fragility to noise and mutations.} The algorithms in this chapter
        are strictly \emph{exact}. Change a single character, delete a comma, or
        insert a single DNA base pair, and KMP, Boyer--Moore, and Rabin--Karp
        report zero matches. In biological sequencing or optical character
        recognition, approximate matching via dynamic programming (edit distance
        and local alignment) becomes necessary.
  \item \textbf{Preprocessing penalties on tiny inputs.} Preprocessing a pattern
        costs $\Theta(m)$ for KMP and $\Theta(m|\Sigma|)$ for a finite
        automaton. If you are searching for a short word in a tiny snippet of text,
        the naive matcher is often faster in wall-clock time because it pays zero
        initialization overhead.
  \item \textbf{Memory footprint on large alphabets.} The transition table
        $\delta$ of a string-matching automaton requires $\Theta(m|\Sigma|)$
        entries. For 16-bit or 32-bit Unicode alphabets, this table easily swells
        to tens of megabytes, thrashing CPU caches. KMP's prefix function $\pi$
        circumvents this by paying an occasional amortized fall-back loop, independent
        of alphabet size.
  \item \textbf{Adversarial collisions in Rabin--Karp.} Rabin--Karp achieves its
        fast average time only when hash collisions are rare. If an adversary
        knows the chosen prime modulus $q$, they can construct texts where every
        shift produces a spurious collision, degrading the algorithm to its
        $\Theta(nm)$ worst-case crawl.
\end{itemize}

\begin{recoverybox}
If only one sentence survives from this chapter, let it be this: the fast string
matchers all win by \emph{never re-reading the text} -- they invest a little
preprocessing of the pattern so that, when a match breaks down, they already know
exactly how far to slide without going back. Rabin--Karp buys that with a
fingerprint; the automaton with a table; KMP with the prefix function.
\end{recoverybox}

In this chapter, we tracked exact patterns across one-dimensional sequences of discrete characters. In Chapter~32, we step from discrete strings to continuous space: computational geometry, where objects are points, segments, and polygons in the plane, and the fundamental questions concern spatial orientation, convex hulls, and geometric intersections.

\section{Exercises}
\label{ch31:exercises}

\begin{exercise}[Warmup: trace the naive matcher]
Run \textsc{Naive-String-Matcher} on the text $T=\texttt{aaabaaaab}$ and pattern
$P=\texttt{aaab}$. List every shift, the number of character comparisons it
makes at that shift, and which shifts are valid. What is the total comparison
count?
\end{exercise}

\begin{exercise}[The quadratic worst case]
Exhibit a text and pattern, each parameterized by $n$ and $m$, on which the naive
matcher makes exactly $(n-m+1)\,m$ character comparisons. Argue briefly why no
shift can be skipped on your input.
\end{exercise}

\begin{exercise}[Rolling hash by hand]
With radix $d=10$ and prime $q=11$, compute the Rabin--Karp fingerprint of the
pattern $P=\texttt{2718}$ (so $m=4$), and the fingerprint of the first
length-$4$ window $\texttt{9271}$ of $T=\texttt{927183}$. Then use the rolling
update to obtain the fingerprint of the second window $\texttt{2718}$ and the
third window $\texttt{7183}$. Show every modular step, and say which window is a
real match for $P$.
\end{exercise}

\begin{exercise}[Spurious hits]
For $d=10$ and $q=13$, find two distinct $3$-digit strings that have the same
Rabin--Karp fingerprint. Explain why the verification step in
\textsc{Rabin-Karp-Matcher} is not optional, using your example.
\end{exercise}

\begin{exercise}[Choosing the prime]
Suppose your text has length $n=10^6$ and you want the \emph{expected} number of
spurious hits to be at most $1$. Using the estimate that the expected number of
spurious hits is about $n/q$, what is the smallest prime $q$ that suffices? How
large does the machine word need to be so that $d\cdot q$ fits, for byte data
($d=256$)?
\end{exercise}

\begin{exercise}[Compute a prefix function]
Compute the prefix function $\pi$ for $P=\texttt{aabaabaaa}$ by hand, showing the
longest prefix-that-is-a-suffix at each position. Then verify your answer by
running \textsc{Compute-Prefix-Function} and listing the value of $k$ at the end
of each iteration of the outer loop.
\end{exercise}

\begin{exercise}[KMP trace]
For the pattern $P=\texttt{ababaca}$ with prefix table
$\pi=[0,0,1,2,3,0,1]$, trace \textsc{KMP-Matcher} on
$T=\texttt{abacababacaba}$. Make a table with columns $i$, $T[i]$, the value of
$q$ before, any fall-backs taken, and $q$ after. Report all valid shifts.
\end{exercise}

\begin{exercise}[Building one automaton transition]
For the pattern $P=\texttt{aabab}$, compute $\delta(4,\texttt{a})$ and
$\delta(4,\texttt{b})$ from the definition (longest prefix of $P$ that is a
suffix of $P_4\,a$). Recall $P_4=\texttt{aaba}$. Show the suffixes you test.
\end{exercise}

\begin{exercise}[Why $q$ never overshoots]
Prove the invariant that drives KMP's correctness: at the top of the outer
\textbf{for} loop in \textsc{KMP-Matcher}, $q$ equals the length of the longest
prefix of $P$ that is a suffix of $T[1\,..\,i-1]$. (Hint: induct on $i$, using
the definition of the prefix function for the mismatch case.)
\end{exercise}

\begin{exercise}[Amortized counting]
Consider the pattern $P=\texttt{aaaaab}$ ($m=6$) and the text of $n$ \texttt{a}'s.
Trace how the state $q$ rises and how the inner \textbf{while} loop fires across
the scan. Confirm by direct count that the \emph{total} number of inner-loop
iterations over the whole text is at most $n$, matching the potential argument,
even though some individual characters trigger several fall-backs.
\end{exercise}

\begin{exercise}[Challenge: two-dimensional matching]
You are given an $n\times n$ grid of characters and an $m\times m$ pattern grid,
and you want every position where the small grid sits inside the big one.
Describe how to combine the row-wise rolling hash of Rabin--Karp with a second
rolling hash down the columns to solve this in expected $O(n^2)$ time after
$O(m^2)$ preprocessing. What is the analogue of a spurious hit here?
\end{exercise}

\begin{exercise}[Challenge: shrinking the automaton table]
The direct automaton construction costs $O(m^2|\Sigma|)$. Explain how the prefix
function $\pi$ lets you compute the transition $\delta(q,a)$ for any state $q$ and
character $a$ on the fly in amortized $O(1)$, so that you never store the full
$\Theta(m|\Sigma|)$ table. (Hint: $\delta(q,a)$ either advances to $q+1$ when
$P[q+1]=a$, or behaves like $\delta(\pi[q],a)$ otherwise -- which is exactly the
inner loop of KMP.)
\end{exercise}
